\section{Method}
\label{sec:method}

The method takes a corpus before and after an ingestion and returns a ranked
list of test questions to run again. Figure~\ref{fig:pipeline} shows the full
pipeline. Every value given in this section is the default in the code.

\begin{figure}[htbp]
  \centering
  \resizebox{\linewidth}{!}{%
  \begin{tikzpicture}[
      node distance=6mm and 6mm,
      box/.style={draw, rounded corners, align=center, font=\small,
                  minimum height=10mm, minimum width=19mm, inner sep=3pt},
      data/.style={box, fill=gray!10},
      proc/.style={box, fill=blue!8},
      arr/.style={-{Stealth[length=2mm]}},
    ]
    % Row 1: build the corpus, then ingest.
    \node[data] (abstracts) {SciFact\\abstracts};
    \node[proc, right=of abstracts] (chunk) {chunk\\$\le 512$ tokens};
    \node[proc, right=of chunk] (embed) {embed\\bge-small};
    \node[data, right=of embed] (before) {corpus\\before $C$};
    \node[proc, right=of before] (ingest) {ingestion\\traps $N$};
    \node[data, right=of ingest] (after) {corpus after\\$C \cup N$};
    % Row 2: the selector sees both corpora; the ground truth re-runs retrieval.
    \node[proc, below=9mm of ingest, xshift=-12mm] (selector)
      {selector: Mapper diff,\\random, k-means,\\nearest neighbours};
    \node[proc, below=9mm of after, xshift=6mm] (retrieval)
      {top-$k$ retrieval,\\$k = 5$, all questions};
    % Row 3: compare.
    \node[data, below=of selector] (ranked) {ranked questions $S_b$\\($\le$ budget $b$)};
    \node[data, below=of retrieval] (broken) {broken\\questions $B$};
    \node[proc, below=6mm of $(ranked.south)!0.5!(broken.south)$] (metric)
      {detection rate\\$|S_b \cap B| / |B|$};

    \draw[arr] (abstracts) -- (chunk);
    \draw[arr] (chunk) -- (embed);
    \draw[arr] (embed) -- (before);
    \draw[arr] (before) -- (ingest);
    \draw[arr] (ingest) -- (after);
    \draw[arr] (before.south) -- (selector.north -| before.south);
    \draw[arr] (after.south) |- ([yshift=3mm]selector.north east) -- (selector.north east);
    \draw[arr] (after.south -| retrieval.north) -- (retrieval.north);
    \draw[arr] (selector) -- (ranked);
    \draw[arr] (retrieval) -- (broken);
    \draw[arr] (ranked.south) |- (metric.west);
    \draw[arr] (broken.south) |- (metric.east);
  \end{tikzpicture}}
  \caption{Pipeline. Selectors see the corpus before and after an ingestion,
  but never run the failure check. The ground truth comes from re-running
  retrieval on every question, and only the evaluation sees it.}
  \label{fig:pipeline}
\end{figure}

\subsection{Corpus, retrieval and failing questions}
\label{sec:method-retrieval}

We use the SciFact test split~\cite{wadden2020scifact} in the BEIR
format~\cite{thakur2021beir}: 5,183 abstracts and 300 questions with 339
relevance judgments. A chunk is the title, a line break and a run of whole
sentences of one abstract, never longer than the model's limit of 512 tokens.
The 455 abstracts that do not fit are cut into the fewest chunks that fit,
with sentences packed evenly. This gives 5,658 chunks. The correct chunks
$R(q)$ of a question are all chunks of its relevant documents, and every
question has at least one.

Chunks and questions are embedded with
\texttt{bge-small-en-v1.5}~\cite{xiao2024cpack} (pinned revision, 384
dimensions, unit vectors). Retrieval returns the $k = 5$ chunks with the
highest cosine similarity, ties going to the earlier chunk. As in
Section~\ref{sec:background}, a question \emph{fails} when
$R(q) \cap \mathrm{top}_5(q, C) = \emptyset$, and an ingestion \emph{breaks}
it when it passes on $C$ and fails on $C \cup N$. On $C$, 66 of the 300
questions (22\%) already fail and cannot be broken. We chose $k = 5$ because
a RAG system passes a handful of chunks to the language model: with $k = 1$,
39\% of the questions already fail, and with $k = 10$ a trap must outrank ten
chunks, so few questions break.

\subsection{Ingestions and ground truth}
\label{sec:method-traps}

An ingestion appends trap chunks $N$ and leaves the old chunks and the
relevance judgments unchanged. Each trap is built from a relevant chunk of a
target question that passes on $C$, in one of three ways. A
\texttt{near\_copy} alters details (on the histone variant Htz1,
``activated'' becomes ``inhibited''). A \texttt{contradiction} reverses a
finding with rule-based edits (``is associated with'' becomes ``is not
associated with''). A \texttt{mix} joins part of a relevant chunk with part
of an unrelated one (IL2RA in type 1 diabetes with the R language). More
examples are in \texttt{report/trap\_examples.md}. The ground truth is the
set of broken questions, found by re-running retrieval on all 300 questions,
which is exactly the work a selector tries to avoid.

\subsection{Selectors}
\label{sec:method-selectors}

A \emph{selector} receives both corpora, the candidate questions and a
budget $b$, and returns at most $b$ distinct questions, most likely broken
first. Because the list is ranked, a smaller budget selects a prefix of a
larger one. A chunk is \emph{new} if its id is not in $C$ or its text
changed; the candidates are the questions with a correct chunk; remaining
ties are broken by question id.

\paragraph{Mapper graph diff.}
The Mapper graph (Section~\ref{sec:mapper}) is built with
KeplerMapper~\cite{vanveen2019kepler}. None of its four settings follows
from the theory:
\begin{itemize}
  \item \emph{Lens}: the first principal component of the embeddings. A
    practical choice: it is linear, so it applies unchanged to new chunks.
  \item \emph{Intervals}: 15. Measured: in a simulated ingestion it changed
    the fewest nodes that hold no new chunk (20, against 28 with 10
    intervals and 41 with 20).
  \item \emph{Overlap}: 50\%. Measured: against 30\%, it roughly halves the
    number of connected components.
  \item \emph{Clustering}: average linkage on cosine distance, cut at 0.27.
    Measured: it leaves 2.4\% of the chunks alone in a node and no node
    holds more than 4\% of the corpus.
\end{itemize}
The paragraph ``From the circle to the corpus'' in Section~\ref{sec:toy}
explains why the theorems do not cover the resulting graph. On SciFact it
has 1,761 nodes and 2,225 edges. The lens and the cover are fitted on $C$
and reused for $C \cup N$, since refitting would move every interval. We
chose KeplerMapper over tda-mapper~\cite{simi2025tdamapper} because its
cover accepts these fixed limits, and because tda-mapper requires
\texttt{numpy<2}, which conflicts with the rest of the environment.

Each node $v$ after the ingestion is matched to the node $u$ before it that
maximises the Jaccard index with the old chunks of $v$. A node is
\emph{affected} if (1) it contains a new chunk, (2) it has an edge whose
matched nodes had none before, (3) it is one of several nodes matched to the
same node before (a split), or (4) it is the match of several nodes before
(a merge). Each rule can be switched off. The \emph{marked} nodes are the
affected nodes and their neighbours. A question is selected if a correct
chunk lies in a marked node, ranked first by whether that node is affected
or only a neighbour, then by its number of new chunks.

\paragraph{Baselines.}
Three baselines show how much of the result comes from the graph.
\emph{Random} shuffles the candidates with a fixed seed; in expectation it
finds a fraction $b/300$ of the broken questions. \emph{k-means} fits
scikit-learn k-means (10 restarts) on the embeddings of $C$, assigns the
chunks of $C \cup N$ to the nearest centroid and ranks questions by the new
chunks in the cluster of their correct chunk: Mapper's key without edges,
neighbours, splits or merges. It uses as many clusters as Mapper has nodes,
a heuristic for a fair comparison rather than a choice the theory suggests.
\emph{Nearest neighbours} takes the $m = 10$ old chunks most similar to each
new chunk and ranks the questions whose correct chunk is among them by that
similarity. We set $m = 2k$ as a margin, because the neighbours of a chunk
are not exactly those of a question. A larger $m$ never drops a candidate,
but it can add some that rank higher, so the selection for a fixed budget
depends on $m$. Random always fills the budget; the other selectors return
only the questions their rule marks.

\subsection{Evaluation}
\label{sec:method-evaluation}

For an ingestion with broken questions $B \neq \emptyset$ and a selection
$S_b$, the \emph{detection rate} is
$\mathrm{DR}(b) = |S_b \cap B| / |B|$. The budgets are 5\%, 10\% and 20\% of
the questions (15, 30 and 60), and every selector receives all 300 questions as candidates.

There are 30 ingestions: three trap types and seeds 0 to 9. Each seed
samples 20 target questions that pass on $C$ and writes one trap per target.
The three types of a seed share the targets, so they are directly
comparable. k-means is fitted with the seed of the ingestion.

Few traps push a correct chunk out of the top~5, so $B$ is small or empty,
and with one or two broken questions $\mathrm{DR}$ can only be 0, 0.5 or 1.
For each selector, trap type and budget we therefore report the mean and
standard deviation of $\mathrm{DR}$ over the ingestions with
$B \neq \emptyset$, the \emph{pooled} rate $\sum |S_b \cap B| / \sum |B|$,
and the \emph{target rate}, the fraction of the 20 targets in $S_b$, which
gives a denser signal. We also report the \emph{tie block} at the cut, the
number of questions sharing the key of the last one selected: when it is
large, the question-id tie-break decides part of the selection. Mapper is
also run without the new-edge rule and the neighbour expansion, to measure
what the graph structure adds over its clusters.

The stability study repeats the Mapper selector over 18 settings (PCA or
centroid lens; 10, 15 or 20 intervals; overlap 0.3, 0.5 or 0.7; cut fixed at
0.27) and compares each selection with the default's. The main experiment
takes 1,420\,s (24 minutes); the stability study about 17 minutes on
9 processes.
